\section{Cotangent tensors and the permitted boundary}
\label{sec:h:slope}\label{sec:slope}

The induction below produces effective divisors of the form $M+mL$,
where $M$ is a line bundle mapping nontrivially into a cotangent tensor
power.  We must
control the contribution of $M$ using the same boundary that is allowed
in the eventual section.  The finite adjoint volume gives precisely
this control.  The first step is analytic: a relative Bergman metric
remains positive after removing the vanishing of the differential along
multiple fibres.  The second step applies this observation to the
positive part of the tangent sheaf's slope filtration.

Throughout this section, $Z$ is a smooth connected projective complex
variety, $D\geq0$ is an integral divisor, and
\[
                        J=-K_Z+D.
\]
We assume that $J$ has a singular Hermitian metric $h_J$ of semipositive
curvature and that
\begin{equation}\label{eq:h:slope:volume}
              \int_Z |s_D|_{h_J}^2<\infty.
\end{equation}
Here $s_D$, considered as a section of $K_Z+J$, is a $J$-valued top
form, and the norm denotes its associated positive density.  A local
weight $\phi$ for $h_J$ satisfies $|e|_{h_J}^2=e^{-\phi}$ in a
holomorphic frame $e$.  It is plurisubharmonic, so the frame norm has a
strictly positive lower bound on each relatively compact chart.  No
upper bound for that norm is assumed.

\subsection{Removing the ramification from the Bergman metric}

A saturated foliation $\mathcal F\subset T_Z$ is a coherent subsheaf
with torsion-free quotient, closed under the Lie bracket.  Write
\[
 Q=T_Z/\mathcal F,\qquad
 P=\det Q^*=(\bigwedge^{\operatorname{rank}Q}Q^*)^{**}.
\]
Since $Z$ is smooth, $P$ is a line bundle.  If the foliation is induced
by a rational fibration, then $P$ agrees generically with the pullback
of the base canonical bundle.  They can differ along a multiple fibre;
that difference is the point of the following lemma.

\begin{lemma}\label{h:slope:foliation}
Under the assumptions above, let $\mathcal F\subset T_Z$ be an
algebraically integrable saturated foliation with
$0<\operatorname{rank}\mathcal F<\dim Z$.  Then the line bundle
\[
                    D-P=D+\det Q
\]
admits a singular Hermitian metric of semipositive curvature.  In
particular, $c_1(Q)+[D]$ is pseudoeffective.
\end{lemma}

\begin{proof}
Set $n=\dim Z$ and $b=\operatorname{rank}Q$.  Resolve the rational
fibration using Lemma~\ref{lem:prepared-model}.  Thus there are
morphisms
\[
 r:Y\longrightarrow\bar Y,\qquad
 \bar\pi:\bar Y\longrightarrow Z,\qquad
 \bar f:\bar Y\longrightarrow S,
 \qquad \pi=\bar\pi\circ r,\quad f=\bar f\circ r,
\]
where $Y,S$ are smooth projective, $\bar Y$ is normal, $\bar\pi$ and
$r$ are birational, and $\bar f$ is equidimensional with connected
general fibres.  Their dimension is $n-b$.  The foliation is the
saturated relative tangent sheaf at the generic point.

Pull back $s_D$ as a top form with coefficients, and denote the result
by
\[
               \eta\in H^0(Y,K_Y+\pi^*J).
\]
Its divisor is $K_{Y/Z}+\pi^*D$; the relative canonical divisor is the
effective birational Jacobian.  Change of variables off the exceptional
locus shows that the density of $\eta$ for $\pi^*h_J$ has finite total
mass.  This assertion concerns almost-everywhere densities, so no
pointwise product on the metric's polar set is needed.

Over a coordinate chart in the base, let $ds$ be a nonvanishing frame
of $K_S$.  On a smooth fibre $Y_s$, the restriction of
$\eta/f^*ds$ is a section of $K_{Y_s}+\pi^*J|_{Y_s}$.  It is nonzero
on general fibres, and Fubini's theorem gives finite squared norm on
almost every such fibre.  Berndtsson--P\u{a}un's relative Bergman
kernel theorem therefore supplies a semipositive metric on
\[
                       K_{Y/S}+\pi^*J
\]
\cite[Theorem~0.1]{BerndtssonPaun2008}.  Over the smooth-fibre locus, in a holomorphic frame its weight
$\psi$ is the logarithm of the supremum of the squared absolute
coefficients of fibre sections having squared norm at most one.
Only sections integrable for the restricted twisting metric enter
that supremum.

We now compare this metric with the line bundle in the statement.
Remove from $Z$ a closed subset of codimension at least two so that
$\pi$ is an isomorphism over the remaining open set and the sheaves
$\mathcal F,Q$ are locally free there.  At the generic point of every
divisor of this open set the conormal inclusion is
\[
                         f^*\Omega_S^1\longrightarrow Q^*.
\]
For a horizontal divisor it has full rank without a divisorial zero:
the restriction of $f$ to that divisor is generically submersive in
characteristic zero.  Every other divisor under consideration maps
onto a base divisor.  Indeed its strict transform on $\bar Y$ is a
prime divisor, and equidimensionality excludes a base image of
codimension at least two.

Fix such a vertical divisor $E$.  At its general points there are
coordinates on $Y$ and $S$ in which
\begin{equation}\label{eq:h:slope:power-map}
 s_1=z_1^e,\qquad s_i=z_i\quad(2\leq i\leq b),
 \qquad e=\operatorname{ord}_E(f^*\{s_1=0\}).
\end{equation}
To obtain these coordinates, take a point where the reduced map
$E\to\{s_1=0\}$ is submersive, write $f^*s_1=z_1^e u$ with $u$ a
unit, and absorb a local $e$-th root of $u$ into $z_1$.  Complete the
other pulled-back base coordinates to a coordinate system.  In this
chart the saturated conormal bundle is spanned by
$dz_1,\ldots,dz_b$.  The determinant of its displayed inclusion is
$e z_1^{e-1}$.  Consequently its vanishing order is exactly $e-1$.

This same integer occurs in the Bergman metric.  Choose a frame $v$
of $\pi^*J$ and use
\[
  \xi=(dz_1\wedge\cdots\wedge dz_n)\otimes(f^*ds)^{-1}\otimes v
\]
as the frame of $K_{Y/S}+\pi^*J$.  If $a\xi$ is a fibre section,
then, up to a constant of absolute value one from the order of wedge
factors, its coefficient as an actual fibre top form is
\[
                       \frac{a}{e z_1^{e-1}}.
\]
This distinction between relative and fibre coefficients is also the
one used in the definition of the relative Bergman kernel
\cite[Section~1]{BerndtssonPaun2008}.

Choose two concentric coordinate polydiscs whose closures lie in this
chart.  For each fixed $(z_1,\ldots,z_b)$ with $z_1\ne0$, the larger
polydisc in the free variables $z_{b+1},\ldots,z_n$ lies in the same
smooth nearby fibre.  Its radius can be fixed independently of the
base point.  The twisting frame norm has a uniform positive lower
bound there.  The holomorphic submean inequality on that polydisc,
for points in the smaller one, therefore gives
\[
 \|a\xi\|_{L^2(Y_s)}^2
 \ \geq\ c\,|z_1|^{-2(e-1)}
       \int_{\text{free polydisc}}|a|^2\,d\lambda
 \ \geq\ c'|a(z)|^2|z_1|^{-2(e-1)},
\]
where $c,c'>0$ do not depend on the fibre section or the nearby base
point.  For squared norm at most one, this implies
\begin{equation}\label{eq:h:slope:bergman-order}
                    \psi\leq(e-1)\log|z_1|^2+C.
\end{equation}
The estimate holds on the smooth-fibre locus; any exceptional
almost-everywhere qualifications disappear for the psh representative
by the submean inequality.

On the submersion locus where $\pi$ is an isomorphism we have
$K_{Y/S}+\pi^*J=D-f^*K_S$.  The determinant computation above says
that a frame of $D-P$ maps to $z_1^{e-1}$ times a nonvanishing multiple
of a frame of $D-f^*K_S$.  The induced metric on $D-P$ thus has weight
\[
                         \psi-(e-1)\log|z_1|^2
\]
up to a pluriharmonic frame term.  This weight is psh off $E$ and is
locally bounded above by \eqref{eq:h:slope:bergman-order}; hence it
extends psh across $E$.  At a horizontal divisor there is no
ramification correction, and the original Bergman metric already
extends.  The coordinate argument applies over a dense open subset
of each divisor.  The remaining excluded subset has codimension at
least two in the smooth variety $Z$, across which psh weights extend
uniquely.  The extensions respect changes of frame and therefore
define a semipositive metric on the global line bundle $D-P$.
\end{proof}

\subsection{The slope bound and the face used by induction}

A movable class is an element of the closed movable cone
$\operatorname{Mov}_1(Z)\subset N_1(Z)_{\mathbb R}$.  On a smooth
projective variety this cone is dual to the pseudoeffective divisor
cone \cite{BDPP2013}.  For a torsion-free coherent sheaf $E$ of
positive rank and a nonzero movable class $\alpha$, write
\[
 \mu_\alpha(E)=\frac{c_1(E)\cdot\alpha}{\operatorname{rank}E},
 \qquad \mu_{\min,\alpha}(E),\quad\mu_{\max,\alpha}(E)
\]
for its slope and extremal Harder--Narasimhan slopes.  Equivalently,
the minimum is the infimum of slopes of its nonzero torsion-free
quotients, and the maximum is the supremum of slopes of its nonzero
subsheaves.  We use the Harder--Narasimhan filtration and tensor slope
formula for movable classes, including real boundary classes
\cite[Corollary~2.27 and Theorem~4.2]{GKP2016}.

\begin{theorem}[Finite-volume cotangent bound]
\label{h:slope:cotangent}\label{thm:finite-volume-slope}
Let $Z$ be smooth connected projective over $\C$, let $D\ge0$ be
an integral divisor, and suppose that $J=-K_Z+D$ has a semipositive
singular Hermitian metric $h_J$ with $\int_Z|s_D|_{h_J}^2<\infty$.
If $\dim Z>0$,
then every nonzero real movable class $\alpha$ satisfies
\begin{equation}\label{eq:h:slope:min}
                  \mu_{\min,\alpha}(T_Z)\geq-D\cdot\alpha.
\end{equation}
For every integer $k\geq0$ and every line bundle $M$ with a nonzero
morphism $M\to(\Omega_Z^1)^{\otimes k}$, the line bundle
\begin{equation}\label{eq:h:slope:line}
                             -M+kD
\end{equation}
is pseudoeffective.  More generally, if $E$ is a torsion-free subsheaf
of rank $r>0$ in $(\Omega_Z^1)^{\otimes k}$, then
$-\det E+krD$ is pseudoeffective.
\end{theorem}

\begin{proof}
Assume first that $\dim Z>0$, fix $\alpha$, and set
$\delta=D\cdot\alpha\geq0$.  Semipositivity of $J$ gives
\begin{equation}\label{eq:h:slope:total}
                         c_1(T_Z)\cdot\alpha\geq-\delta.
\end{equation}
Let $\lambda_1>\cdots>\lambda_t$ be the slopes of the
Harder--Narasimhan quotients of $T_Z$, with positive ranks
$r_1,\ldots,r_t$.  Suppose, for a contradiction, that
$\lambda_t<-\delta$.  The total degree of the pieces with
nonpositive slopes then satisfies
\begin{equation}\label{eq:h:slope:negative-degree}
       \sum_{\lambda_i\leq0}r_i\lambda_i
                     \leq r_t\lambda_t<-\delta.
\end{equation}
By \eqref{eq:h:slope:total}, some piece has positive slope.
Consequently the positive part $\mathcal F$ of the filtration is
nonzero and proper.  It is saturated, and
\[
 \mu_{\min,\alpha}(\mathcal F)>0,
 \qquad \mu_{\max,\alpha}(T_Z/\mathcal F)\leq0.
\]

The bracket modulo $\mathcal F$ is an $\mathcal O_Z$-linear morphism
$\bigwedge^2\mathcal F\to T_Z/\mathcal F$.  Since the target is
torsion-free, it factors through
$A=(\bigwedge^2\mathcal F)/\mathrm{torsion}$.  When
$\operatorname{rank}\mathcal F>1$, the tensor slope formula, duality,
and the quotient inequality give
$\mu_{\min,\alpha}(A)\geq2\mu_{\min,\alpha}(\mathcal F)>0$.
The maximum slope of the target is nonpositive, so this morphism
vanishes.  If $\operatorname{rank}\mathcal F=1$, then $A=0$.
Thus $\mathcal F$ is a foliation.  Campana--P\u{a}un's algebraicity
criterion now shows that it is algebraically integrable
\cite[Theorem~1.1]{CampanaPaun2019}.

Lemma~\ref{h:slope:foliation} gives
$c_1(T_Z/\mathcal F)\cdot\alpha\geq-\delta$.  Its left side is
exactly the sum in \eqref{eq:h:slope:negative-degree}, a
contradiction.  This proves \eqref{eq:h:slope:min}.  In particular the
argument covers the cases in which the positive part would be zero
or the whole tangent sheaf: the assumed failure of the inequality
has already excluded both cases.

By duality and tensor-product compatibility,
\[
 \mu_{\max,\alpha}((\Omega_Z^1)^{\otimes k})
                 =-k\mu_{\min,\alpha}(T_Z)
                 \leq kD\cdot\alpha
\]
for $k>0$.  A nonzero map from $M$ is injective, because its target is
torsion-free.  Thus $(kD-M)\cdot\alpha\geq0$ for every movable
$\alpha$, and divisor--curve duality proves
\eqref{eq:h:slope:line}.  For the sheaf $E$, the same inequality
applied to its slope gives
$c_1(E)\cdot\alpha\leq krD\cdot\alpha$, which proves the determinant
assertion.  Determinants are taken reflexively; removing
codimension-two singular loci does not change their first Chern
classes.

If $k=0$, the target is $\mathcal O_Z$.  A nonzero map from $M$
provides a nonzero section of $-M$, and a positive-rank subsheaf of
$\mathcal O_Z$ has rank one with effective negative determinant.
These arguments also cover dimension zero.  In that dimension there
is no nonzero cotangent tensor for $k>0$, and no minimum slope of a
zero tangent sheaf is invoked.
\end{proof}

The coefficient of $D$ in this theorem is sharp even for curves.
Indeed, on a smooth projective curve of genus at least two, take a
nonzero holomorphic one-form $\omega$ and put $D=\operatorname{div}(\omega)$.
Then $J=-K_Z+D$ is trivial and its flat metric satisfies
\eqref{eq:h:slope:volume}.  Taking $M=kK_Z$ gives equality in the
cotangent bound; replacing $kD$ by $cD$ with $c<k$ gives negative
degree and therefore cannot give a pseudoeffective class.

\begin{corollary}[A supporting class for the induction]
\label{h:slope:face}
Under the hypotheses of Theorem~\ref{thm:finite-volume-slope},
assume $\dim Z>0$ and let $L$ be a pseudoeffective line bundle for
which $L+D$ is not big.  There exists a nonzero real movable class
$\alpha$ with
\[
                       L\cdot\alpha=D\cdot\alpha=0.
\]
For this class $\mu_{\min,\alpha}(T_Z)\geq0$.  Suppose moreover that
$M\to(\Omega_Z^1)^{\otimes k}$ is nonzero and that effective divisors
$N_i\sim M+m_iL$ are given.  Then
\[
                  M\cdot\alpha=N_i\cdot\alpha=0
                         \quad\text{for every }i.
\]
Every finite nonnegative linear combination of $D,N_1,N_2,\ldots$
has zero intersection with $\alpha$ and is not big.
\end{corollary}

\begin{proof}
The class of $L+D$ is on the boundary of the closed pseudoeffective
cone, whose interior is the big cone.  A supporting hyperplane and
cone duality give a nonzero real movable class $\alpha$ annihilating
it.  Both $L$ and $D$ pair nonnegatively with $\alpha$, so they pair
with it individually to zero.  Theorem~\ref{h:slope:cotangent} gives
the tangent inequality and $M\cdot\alpha\leq0$.  Effectivity of each
$N_i$ gives the reverse inequality through
$N_i\cdot\alpha=M\cdot\alpha$.  Finally an interior point of a
full-dimensional closed convex cone pairs strictly positively with
every nonzero member of its dual cone; hence no indicated combination
is big.
\end{proof}

The supporting class in this corollary can lie on the boundary of the
movable cone and need not be rational or represented by a single
moving curve.  Its role is numerical: it rules out a generically
finite linear system made from the displayed effective divisors and
thereby forces the dimension drop in the subsequent argument.
